% rel_chapter_5_sec53-54.tex
% Relativistic overstability with rotation + magnetic field (Ch V §53-54)
% Conventions: (-,+,+,+) signature, c kept explicit, u^μ u_μ = -c²
% See RELATIVISTIC_CONVENTIONS.md for notation and macros.

%% ==================================================================
%% Section 53 — Relativistic
%% ==================================================================
\section{The case when instability sets in as overstability --- relativistic}\label{sec:rel-53}

In the classical treatment of \S\,53, overstability in a rotating, magnetised
layer of liquid metal was analysed using the Navier--Stokes--Fourier system,
which is inherently acausal (infinite signal speeds for heat and viscous
stress).  For a relativistic treatment the Israel--Stewart causal dissipation
framework replaces Fourier's law and the Newtonian stress--strain relation
with relaxation equations possessing finite propagation speeds.  The combined
effects of rotation (Taylor number $T$) and magnetic field (Chandrasekhar
number $Q$) must now be augmented by the three relaxation times
$\tau_{\Pi}$, $\taupi$, and $\tauq$ and by causality bounds on the
Alfv\'en and sound speeds.

\subsection{Linearised relativistic perturbation equations}

We consider a horizontally infinite layer of relativistic conducting fluid
confined between $z=0$ and $z=d$, rotating with angular velocity
$\vect{\Omega}=\Omega\,\hat{\vect{z}}$ and permeated by a uniform vertical
magnetic field $\vect{B}_0 = B_0\,\hat{\vect{z}}$.  The equilibrium
4-velocity is $u^{\mu}_0 = \Lf\,(c,\,\vect{v}_{\mathrm{rot}})$ where
$\vect{v}_{\mathrm{rot}} = \vect{\Omega}\times\vect{r}$ and we work in the
co-rotating frame so that inertial effects enter through the Coriolis and
centrifugal terms.

The background energy--momentum tensor is
\begin{equation}
  T^{\mu\nu}_0 = \frac{\varepsilon_0+p_0}{c^{2}}\,u^{\mu}_0\,u^{\nu}_0
  + p_0\,g^{\mu\nu}
  + \frac{1}{2}\,b^{2}\!\left(\frac{u^{\mu}_0 u^{\nu}_0}{c^{2}}
  + \frac{1}{2}\,g^{\mu\nu}\right) - b^{\mu}b^{\nu},
  \label{eq:rel53-T0}
\end{equation}
where $b^{\mu}$ is the magnetic-field 4-vector (cf.\ Chapter~\ref{ch:rel-mhd-framework}).

Linearising the Israel--Stewart equations about this equilibrium, the
perturbation in the heat flux $\delta q^{\mu}$ obeys
\begin{equation}
  \tauq\,u^{\alpha}_0\,\nabla_{\alpha}\,\delta q^{\mu}
  + \delta q^{\mu}
  = -\kappa_T\!\left(\Delta^{\mu\nu}_0\,\nabla_{\nu}\delta T
    + \frac{T_0}{c^{2}}\,u^{\mu}_0\,u^{\alpha}_0\,\nabla_{\alpha}\delta T
  \right),
  \label{eq:rel53-IS-heat}
\end{equation}
and analogously for the shear-stress perturbation $\delta\pi^{\mu\nu}$:
\begin{equation}
  \taupi\,u^{\alpha}_0\,\nabla_{\alpha}\,\delta\pi^{\langle\mu\nu\rangle}
  + \delta\pi^{\mu\nu}
  = 2\,\shearvisc\,\delta\sigma^{\mu\nu}.
  \label{eq:rel53-IS-shear}
\end{equation}
The bulk viscous perturbation $\delta\Pi$ satisfies
\begin{equation}
  \tau_{\Pi}\,u^{\alpha}_0\,\nabla_{\alpha}\,\delta\Pi + \delta\Pi
  = -\bulkvisc\,\nabla_{\mu}\,\delta u^{\mu}.
  \label{eq:rel53-IS-bulk}
\end{equation}

\subsection{Israel--Stewart relaxation effects on oscillatory modes}

The most distinctive feature of the relativistic formulation is that
each dissipative channel introduces a relaxation time.  In the frequency
domain, replacing $u^{\alpha}_0\nabla_{\alpha}\to -i\omega$ in the
co-rotating frame, the Israel--Stewart closures \eqref{eq:rel53-IS-heat}---%
\eqref{eq:rel53-IS-bulk} introduce frequency-dependent effective
transport coefficients:
\begin{equation}
  \kappa_T^{\mathrm{eff}}(\omega)
  = \frac{\kappa_T}{1 - i\omega\tauq},\qquad
  \shearvisc^{\mathrm{eff}}(\omega)
  = \frac{\shearvisc}{1 - i\omega\taupi},\qquad
  \bulkvisc^{\mathrm{eff}}(\omega)
  = \frac{\bulkvisc}{1 - i\omega\tau_{\Pi}}.
  \label{eq:rel53-eff-transport}
\end{equation}
These reduce to the classical (Navier--Stokes--Fourier) values when
$\omega\tau\to 0$ and impose an upper cutoff $\omega_{\max}\sim 1/\tau$
beyond which dissipation is suppressed --- precisely the mechanism that
restores causality.

For overstable modes of frequency $\omega = \omega_r + i\omega_i$ (with
$\omega_i>0$ the growth rate), the classical analysis of \S\,53 had
$\sigma = i\sigma_i$ (purely imaginary in the marginal state).  In the
relativistic theory the marginal condition $\omega_i = 0$ still picks out
neutral oscillatory modes, but their real frequencies are modified:
\begin{equation}
  \omega_r^{(\mathrm{rel})} = \omega_r^{(\mathrm{class})}
  \left[1 + \omega_r^{2}\,\tauq^{2}\,
  \frac{p_1\,(1-p_1)}{(1+\omega_r^{2}\tauq^{2})}\right]^{1/2}
  + \mathcal{O}(v^{2}/c^{2}).
  \label{eq:rel53-freq-shift}
\end{equation}
Since $\tauq>0$, the effect is always to \emph{increase} the oscillation
frequency relative to the classical prediction --- the causal heat-flux
propagation stiffens the restoring mechanism.

\subsection{Relativistic characteristic equation}

Adopting the same free-boundary geometry ($W = W_0\sin\pi z$, non-conducting
exterior), the relativistic analogue of equation~(62) of \S\,53 is
\begin{align}
  R_*^{(\mathrm{rel})}\,a^{2}\,&(n^{2}{+}a^{2}{+}p_{2}^{(\mathrm{eff})}\sigma)
  \bigl[(n^{2}{+}a^{2}{+}\sigma^{(\mathrm{eff})})(n^{2}{+}a^{2}{+}p_{2}^{(\mathrm{eff})}\sigma)
  {+}Q_{\mathrm{rel}}\,n^{2}\bigr]
  \notag\\
  &= (n^{2}{+}a^{2}{+}p_{1}^{(\mathrm{eff})}\sigma)
  \Bigl\{(n^{2}{+}a^{2})
  \bigl[(n^{2}{+}a^{2}{+}\sigma^{(\mathrm{eff})})(n^{2}{+}a^{2}{+}p_{2}^{(\mathrm{eff})}\sigma)
  {+}Q_{\mathrm{rel}}\,n^{2}\bigr]^{2}
  \notag\\
  &\qquad {}+ T_{\mathrm{rel}}\,n^{2}\,(n^{2}{+}a^{2}{+}p_{2}^{(\mathrm{eff})}\sigma)^{2}
  \Bigr\},
  \label{eq:rel53-chareq}
\end{align}
where the relativistic Prandtl-like ratios carry relaxation corrections,
\begin{equation}
  p_{1}^{(\mathrm{eff})} = \frac{\shearvisc^{\mathrm{eff}}}
  {\kappa_T^{\mathrm{eff}}}
  = p_1\,\frac{1-i\omega\tauq}{1-i\omega\taupi},\qquad
  p_{2}^{(\mathrm{eff})} = \frac{\shearvisc^{\mathrm{eff}}}
  {\eta^{\mathrm{eff}}}
  = p_2\,\frac{1-i\omega\tau_{\eta}}{1-i\omega\taupi},
  \label{eq:rel53-peff}
\end{equation}
and the dimensionless numbers acquire relativistic corrections:
\begin{align}
  R_*^{(\mathrm{rel})} &= R_*\!\left(1 + \mathcal{C}_R\,\frac{v^{2}}{c^{2}}\right),
  \label{eq:rel53-Rrel}\\[4pt]
  T_{\mathrm{rel}} &= T\!\left(1 + \mathcal{C}_T\,\frac{\Omega^{2}d^{2}}{c^{2}}\right),
  \label{eq:rel53-Trel}\\[4pt]
  Q_{\mathrm{rel}} &= Q\!\left(1 + \mathcal{C}_Q\,\frac{\vA^{2}}{c^{2}}\right),
  \label{eq:rel53-Qrel}
\end{align}
with $\mathcal{C}_R$, $\mathcal{C}_T$, $\mathcal{C}_Q$ dimensionless
coefficients of order unity determined by the equation of state.

\subsection{Approximate solution for liquid metals (relativistic)}\label{subsec:rel53-approx}

For liquid metals such as mercury, $p_2\ll p_1\ll 1$.  In the non-relativistic
regime (\S\,53(a)), all terms proportional to $p_2$ were dropped.  In the
relativistic extension, the magnetic relaxation time $\tau_{\eta}$ is
likewise negligible (it scales with the resistive diffusion time, far shorter
than other hydrodynamic scales), so $p_2^{(\mathrm{eff})}\to p_2$.  However
the \emph{thermal} relaxation time $\tauq$ can be non-negligible even for
liquid metals when the oscillation frequency $\omega_r$ is high enough.

Suppressing $p_2$ as in the classical case, the relativistic analogue of
equation~\eqref{eq:5-70} for $\sigma_i^{2}$ becomes
\begin{equation}
  \sigma_i^{2} = \frac{T_{\mathrm{rel}}}{1{+}x}\;
  \frac{(1{+}x)^{2}(1{-}p_1^{(\mathrm{eff})}) - p_1^{(\mathrm{eff})}\,Q_{\mathrm{rel}}}
  {(1{+}x)(1{+}p_1^{(\mathrm{eff})})+p_1^{(\mathrm{eff})}\,Q_{\mathrm{rel}}}
  - \left[\frac{(1{+}x)^{2}{+}Q_{\mathrm{rel}}}{1{+}x}\right]^{2}\!,
  \label{eq:rel53-sigmai}
\end{equation}
and the Rayleigh number is
\begin{equation}
  R_*^{(\mathrm{rel})} = 2\,\frac{1+x}{x}\;\frac{(1{+}x)^{2}{+}Q_{\mathrm{rel}}}
  {(1{+}x)^{2}(1{-}p_1^{(\mathrm{eff})})}\;
  \bigl[(1{+}x)^{2}+p_1^{(\mathrm{eff})}\,\sigma_i^{2}\bigr].
  \label{eq:rel53-Rstar}
\end{equation}
In the formal limit $c\to\infty$ (whence $\tauq\to 0$, $\vA/c\to 0$,
$\Omega d/c\to 0$) these reduce identically to equations~\eqref{eq:5-70}
and~\eqref{eq:5-71}.

\subsection{Mercury results --- relativistic corrections}\label{subsec:rel53-mercury}

For mercury at room temperature, the characteristic thermal relaxation
time may be estimated from kinetic theory as
\begin{equation}
  \tauq \approx \frac{3\,\kappa_T}{c_s^{2}}
  \sim 10^{-12}\;\mathrm{s},
  \label{eq:rel53-tauq-Hg}
\end{equation}
where $c_s\approx 1{,}450\;\mathrm{m\,s^{-1}}$ is the sound speed in mercury.
Since typical overstable oscillation frequencies in Nakagawa's experiments
are $\omega_r\lesssim 10^{2}\;\mathrm{s^{-1}}$ (see Table~\ref{tab:XIX}),
we have $\omega_r\,\tauq \sim 10^{-10}\ll 1$.

Similarly, the Alfv\'en speed for fields $B_0\leq 5{,}000\;\mathrm{G}$ in
mercury ($\rho\approx 1.36\times10^{4}\;\mathrm{kg\,m^{-3}}$) is
\begin{equation}
  \vA = \frac{B_0}{\sqrt{\mu_0\,\rho}}
  \lesssim 1.2\;\mathrm{m\,s^{-1}}
  \quad\Longrightarrow\quad
  \frac{\vA^{2}}{c^{2}} \sim 10^{-17},
  \label{eq:rel53-vA-Hg}
\end{equation}
and the rotational velocity at $\Omega = 5\;\mathrm{rev/min}$,
$r\sim 12\;\mathrm{cm}$ is $v_{\mathrm{rot}}\sim 0.06\;\mathrm{m\,s^{-1}}$,
giving $v_{\mathrm{rot}}^{2}/c^{2}\sim 10^{-20}$.

Consequently, for mercury under laboratory conditions, all relativistic
corrections --- whether from Israel--Stewart relaxation, finite
$\vA/c$, or rotational frame-dragging --- are utterly negligible:
\begin{equation}
  \frac{R_c^{(\mathrm{rel})} - R_c^{(\mathrm{class})}}{R_c^{(\mathrm{class})}}
  \sim \max\!\left(\omega_r^{2}\tauq^{2},\;
  \frac{\vA^{2}}{c^{2}},\;
  \frac{\Omega^{2}d^{2}}{c^{2}}\right)
  \lesssim 10^{-17}.
  \label{eq:rel53-Hg-correction}
\end{equation}
The classical results of Table~\ref{tab:XIX} therefore remain accurate
to all experimentally accessible precision.  This confirms that
Chandrasekhar's original overstability analysis is the correct
non-relativistic limit of the causal theory.

\subsection{Regimes where relativistic corrections matter}

Relativistic corrections become significant when any of the following
conditions hold:
\begin{enumerate}
  \item \textbf{Hot, dense plasmas} (e.g.\ neutron-star envelopes):
    the sound speed $c_s$ approaches $c/\!\sqrt{3}$, and the thermal
    relaxation time $\tauq$ is comparable to the oscillation period.
  \item \textbf{Strongly magnetised plasmas} (magnetar interiors):
    $\vA/c$ can approach unity when $B\gtrsim 10^{15}\;\mathrm{G}$,
    making $Q_{\mathrm{rel}}$ substantially different from $Q$.
  \item \textbf{Millisecond pulsars and accretion disks}:
    $\Omega\,d/c$ can be non-negligible, and frame-dragging corrections
    to $T_{\mathrm{rel}}$ become relevant.
\end{enumerate}
In all these cases the full Israel--Stewart treatment and the causality-bounded
transport coefficients of equations~\eqref{eq:rel53-eff-transport} are
essential for a self-consistent analysis.


%% ==================================================================
%% Section 54 — Relativistic
%% ==================================================================
\section{Experiments and astrophysical applications --- relativistic}\label{sec:rel-54}

The experiments of Nakagawa described in \S\,54 were carried out with
mercury under laboratory conditions where, as shown in
\S\,\ref{subsec:rel53-mercury}, relativistic corrections are negligible.
No laboratory experiment to date has probed the relativistic regime of
combined rotation and magnetic field in convective instability.  The
natural arena for these effects is astrophysics.

\subsection{Combined rotation and magnetic field: magnetar crusts}

Magnetar crusts provide perhaps the most extreme environment in which the
interplay of rotation, magnetic field, and thermal convection operates in
the relativistic regime.  Surface magnetic fields of order
$B\sim 10^{14}$--$10^{15}\;\mathrm{G}$ yield
\begin{equation}
  \frac{\vA^{2}}{c^{2}}
  = \frac{b^{2}}{4\pi\,\enthalpy + b^{2}}
  \sim 10^{-2}\text{--}10^{-1},
  \label{eq:rel54-vA-magnetar}
\end{equation}
where $\enthalpy = (\varepsilon+p)/c^{2}$ is the relativistic enthalpy
density of the crustal material.  At these field strengths:
\begin{itemize}
  \item The Chandrasekhar number $Q_{\mathrm{rel}}$ departs significantly
    from its Newtonian value; the $\vA^{2}/c^{2}$ correction in
    equation~\eqref{eq:rel53-Qrel} contributes at the $1$--$10\%$ level.
  \item The fast magnetosonic speed replaces the Alfv\'en speed as the
    relevant signal velocity for the magnetic tension restoring force:
    \begin{equation}
      v_f^{2} = \cs^{2} + \vA^{2} - \frac{\cs^{2}\,\vA^{2}}{c^{2}}
      < c^{2},
      \label{eq:rel54-vfast}
    \end{equation}
    ensuring causality is respected.
  \item The spin periods of magnetars ($P\sim 2$--$12\;\mathrm{s}$)
    correspond to $\Omega\,R_{\star}/c\sim 10^{-4}$, so rotational
    relativistic corrections remain modest.  However, frame-dragging
    (Lense--Thirring) effects in the interior can shift the effective
    Taylor number:
    \begin{equation}
      T_{\mathrm{rel}} = T\!\left(1
      + \frac{4\,G\,I\,\Omega}{c^{2}\,R_{\star}^{3}}\right),
      \label{eq:rel54-T-framedrag}
    \end{equation}
    where $I$ is the stellar moment of inertia and $R_{\star}$ the stellar
    radius.
\end{itemize}

The overstability criterion of \S\,\ref{sec:rel-53} implies that in
magnetar crusts, oscillatory convective modes may be excited at Rayleigh
numbers \emph{below} those required for stationary convection.  These
oscillatory modes have been proposed as a mechanism for quasi-periodic
oscillations (QPOs) observed in the tails of giant magnetar flares
\citep{Israel2005,Strohmayer2005,Watts2006}.

\subsection{Magnetised accretion disks}

In geometrically thick accretion disks around black holes, the combined
effects of differential rotation, magnetic field, and thermal stratification
produce a rich landscape of instabilities.  The classical analysis of \S\S\,51--53
provides the foundation for understanding how the magneto-rotational
instability (MRI) interacts with convection when a vertical temperature
gradient is maintained.

In the relativistic regime, three new features arise:
\begin{enumerate}
  \item \textbf{Relativistic enthalpy coupling.}  The inertia of the fluid
    is determined by $\enthalpy = (\varepsilon+p)/c^{2}$ rather than
    $\rho$ alone.  For radiation-dominated disks near the innermost
    stable circular orbit (ISCO), $p\sim\varepsilon/3$, so $\enthalpy$
    exceeds $\rho$ by a factor $\sim 4/3$.  This modifies both the
    Alfv\'en speed and the Rayleigh number.

  \item \textbf{Causal heat transport.}  In radiation-dominated disks
    the photon mean free path sets $\tauq \sim \ell_{\gamma}/c$.  For
    optically thick disks, $\tauq$ can be comparable to the local
    dynamical time $\Omega^{-1}$, making the Israel--Stewart relaxation
    effects on oscillatory modes (\S\,\ref{sec:rel-53}) quantitatively
    important.

  \item \textbf{Frame-dragging.}  Near a spinning (Kerr) black hole with
    dimensionless spin $a_{\star}$, the Lense--Thirring angular velocity
    $\omega_{\mathrm{LT}}\sim G M a_{\star}/(c^{2}\,r^{2})$ modifies the
    effective rotation rate entering the Taylor number.  This effect is
    strongest near the ISCO, where the convective and MRI growth rates are
    also largest.
\end{enumerate}

The interplay between convection and MRI in such environments is an
active area of research; general-relativistic radiation-MHD simulations
(e.g.\ \citealt{McKinney2014,Sadowski2015}) have begun to capture these
effects, though a clean separation into the individual instability
mechanisms of Chandrasekhar's classical theory remains a theoretical
challenge.

\subsection{Neutron-star oceans}

The outermost layers of accreting neutron stars form a thin ``ocean'' of
nuclear-burning material.  At densities $\rho\sim 10^{6}$--$10^{9}\;\mathrm{g\,cm^{-3}}$
and temperatures $T\sim 10^{8}\;\mathrm{K}$, the sound speed is
$\cs\sim 10^{8}\;\mathrm{cm\,s^{-1}}$ (i.e.\ $\cs/c\sim 3\times10^{-3}$),
and surface fields can reach $10^{8}$--$10^{12}\;\mathrm{G}$.
The spin frequencies of accreting millisecond pulsars
($\nu_{\mathrm{spin}}\sim 300$--$700\;\mathrm{Hz}$)
give $\Omega\,d/c\sim 10^{-6}$--$10^{-4}$ for ocean depth $d\sim 10\;\mathrm{m}$.

In these systems the dominant relativistic correction comes not from kinematics
but from the equation of state: the partially degenerate, strongly
interacting electron gas modifies the Prandtl number $p_1$ and the effective
thermal diffusivity.  The Israel--Stewart relaxation time $\tauq$ for
electron--ion collisions is $\sim 10^{-15}\;\mathrm{s}$, again negligibly
short compared with convective time-scales.  Nevertheless, the surface
gravity $g\sim 2\times10^{14}\;\mathrm{cm\,s^{-2}}$ is large enough that
the gravitational redshift correction
\begin{equation}
  R^{(\mathrm{rel})} = R\!\left(1 - \frac{2GM}{c^{2}\,R_{\star}}\right)^{-1}
  \label{eq:rel54-redshift-R}
\end{equation}
amounts to a $\sim 30\%$ increase in the effective Rayleigh number for
a typical neutron star.  This general-relativistic correction to the
hydrostatic background is therefore the leading ``relativistic'' effect
in neutron-star ocean convection, rather than the kinematic or
Israel--Stewart corrections emphasised elsewhere.


%% ==================================================================
%% BIBLIOGRAPHICAL NOTES — Relativistic
%% ==================================================================
\section*{BIBLIOGRAPHICAL NOTES --- Relativistic extensions of \S\S\,53--54}
\addcontentsline{toc}{section}{Bibliographical Notes --- Relativistic \S\S\,53--54}

\paragraph{Israel--Stewart theory.}
The causal thermodynamics framework used throughout this section is due to:
\begin{enumerate}
  \item W.~\textsc{Israel} and J.\,M.~\textsc{Stewart}, `Transient relativistic
    thermodynamics and kinetic theory', \textit{Ann.\ Phys.\ (N.Y.)}\ \textbf{118},
    341--72 (1979).
  \item W.~\textsc{Israel}, `Nonstationary irreversible thermodynamics: a causal
    relativistic theory', \textit{ibid.}\ \textbf{100}, 310--31 (1976).
  \item W.\,A.~\textsc{Hiscock} and L.~\textsc{Lindblom}, `Stability and causality
    in dissipative relativistic fluids', \textit{Ann.\ Phys.\ (N.Y.)}\ \textbf{151},
    466--96 (1983).
\end{enumerate}
For a modern review of relativistic dissipative hydrodynamics:
\begin{enumerate}\setcounter{enumi}{3}
  \item P.~\textsc{Romatschke} and U.~\textsc{Romatschke},
    \textit{Relativistic Fluid Dynamics In and Out of Equilibrium}
    (Cambridge University Press, 2019).
\end{enumerate}

\paragraph{Relativistic MHD stability.}
The extension of Chandrasekhar's magneto-convective analysis to the
relativistic domain draws on:
\begin{enumerate}\setcounter{enumi}{4}
  \item A.\,M.~\textsc{Anile}, \textit{Relativistic Fluids and Magneto-Fluids}
    (Cambridge University Press, 1989).
  \item A.~\textsc{Lichnerowicz}, \textit{Relativistic Hydrodynamics and
    Magnetohydrodynamics} (W.\,A.\ Benjamin, New York, 1967).
\end{enumerate}

\paragraph{Magnetar QPOs and crustal oscillations.}
\begin{enumerate}\setcounter{enumi}{6}
  \item G.\,L.~\textsc{Israel} \textit{et al.}, `The discovery of rapid X-ray
    oscillations in the tail of the SGR~1806$-$20 hyperflare',
    \textit{Astrophys.\ J.}\ \textbf{628}, L53--L56 (2005).
  \item T.\,E.~\textsc{Strohmayer} and A.\,L.~\textsc{Watts}, `Discovery of fast
    X-ray oscillations during the 1998 giant flare from SGR~1900+14',
    \textit{Astrophys.\ J.}\ \textbf{632}, L111--L114 (2005).
  \item A.\,L.~\textsc{Watts} and T.\,E.~\textsc{Strohmayer}, `Detection with RHESSI
    of high-frequency X-ray oscillations in the tailof the 2004~December~27
    hyperflare from SGR~1806$-$20',
    \textit{Astrophys.\ J.}\ \textbf{637}, L117--L120 (2006).
  \item N.~\textsc{Andersson}, K.\,D.~\textsc{Kokkotas}, and B.\,F.~\textsc{Schutz},
    `Gravitational radiation limit on the spin of young neutron stars',
    \textit{Astrophys.\ J.}\ \textbf{510}, 846--53 (1999).
\end{enumerate}

\paragraph{Relativistic convection in accretion disks.}
\begin{enumerate}\setcounter{enumi}{10}
  \item J.\,C.~\textsc{McKinney} \textit{et al.}, `Three-dimensional general
    relativistic radiation magnetohydrodynamical simulation of super-Eddington
    accretion, using a new code HARMRAD with M1 closure',
    \textit{Mon.\ Not.\ R.\ Astron.\ Soc.}\ \textbf{441}, 3177--208 (2014).
  \item A.~\textsc{Sa\c{d}owski} \textit{et al.}, `Magnetohydrodynamical
    simulations of black hole accretion on horizon scales',
    \textit{Mon.\ Not.\ R.\ Astron.\ Soc.}\ \textbf{447}, 49--71 (2015).
\end{enumerate}

\paragraph{Neutron-star ocean convection.}
\begin{enumerate}\setcounter{enumi}{12}
  \item A.~\textsc{Cumming} and L.~\textsc{Bildsten}, `Carbon flashes in the heavy
    element ocean on accreting neutron stars',
    \textit{Astrophys.\ J.}\ \textbf{559}, L127--L130 (2001).
  \item H.\,C.~\textsc{Spruit}, `Convective collapse of white dwarf envelopes',
    \textit{Astron.\ Astrophys.}\ \textbf{108}, 348--55 (1982).
  \item E.\,F.~\textsc{Brown}, `Nuclear heating and melted layers in the inner
    crust of an accreting neutron star',
    \textit{Astrophys.\ J.}\ \textbf{531}, 988--1002 (2000).
\end{enumerate}
